\section{Personnel Management}
{\color{sub}\footnotesize 8 hrs \gap$\cdot$\gap asked in \mk{28 of 29 papers}
\gap$\cdot$\gap 7 syllabus topics \gap$\cdot$\gap \mk{400 marks, 17.0\% of the archive}
$\rightarrow$ about 9 of every 80 marks, nearly always one full 8-mark question.}

\vspace{3pt}\noindent
\colorbox{gdL}{\makebox[\dimexpr\textwidth-2\fboxsep][l]{%
  \textcolor{gd}{\bfseries\footnotesize READ THIS FIRST}}}
\par\nobreak\vspace{3pt}

Chapter 2 follows the life of an employee: \textbf{plan} how many people are needed
$\rightarrow$ \textbf{analyse} the job $\rightarrow$ \textbf{recruit and select}
$\rightarrow$ \textbf{train} $\rightarrow$ \textbf{appraise} $\rightarrow$ \textbf{pay}.
Every question picks two stops on that line.
\begin{enumerate}[label=\arabic*., leftmargin=1.4em, itemsep=1pt]
  \item \textbf{Four lists carry the chapter}: definition + functions of personnel management
    (17 papers), recruitment and selection steps (13), factors affecting wage and salary
    structure (12), job analysis with job description and specification (10).
  \item \textbf{Personnel policy} (7 papers) is always asked from the new employee's side:
    ``why discuss the handbook at hiring?''
  \item \textbf{No numerical has ever been set}, but ``how are wages calculated?'' (75 Ch) is
    answered by the incentive-plan formulas in $\S$2.7, each with one worked example.
\end{enumerate}
Tier chips count \textbf{papers out of 29}: \tS{n} 8+ \gap \tF{n} 4--7 \gap \tP{n} 1--3
inside an 8-mark question \gap \tL{n} 1--3, short note only. Bold year = Regular exam.

\hr

%=======================================================================
\T{2.1 Personnel Management: Meaning, Functions, Role}

\Q{Define personnel management. Explain its functions and its role in the organization}{\m{1}\m{2}\m{3}\m{4}\m{5}\m{8}}
{\color{sub}\footnotesize \tS{17} \yr{\textbf{82 Bh}, 81 Ba, \textbf{80 Bh}, 80 Ba, \textbf{76 Ch}, 76 Ash, \textbf{75 Ch}, 75 Ash, 74 Ash, 72 Ka, \textbf{71 Ch}, 71 Shr, \textbf{70 Ch}, 70 Asa, \textbf{68 Ch}, \textbf{68 Ba}, \textbf{65 Shr}} \gap$\cdot$\gap \mk{the most asked question of the chapter}}

\begin{itemize}
  \item \textbf{Flippo}: ``the planning, organizing, directing and controlling of the
    \textbf{procurement, development, compensation, integration and maintenance} of personnel,
    to accomplish organizational, individual and social goals.''
  \item Simple: the part of management that \textbf{obtains, develops, uses and keeps a
    satisfied workforce}. Also called human resource management.
  \item \textbf{Aims}: best out of people; maximum individual development; better service to
    society; good labour relations.
  \item \textbf{Three aspects} (AJ Sir): \textbf{welfare} (working conditions, canteen,
    housing), \textbf{labour / personnel} (recruitment, pay, promotion),
    \textbf{industrial relations} (unions, disputes, collective bargaining).
\end{itemize}

\sbs{%
\lead{Managerial functions}
\begin{itemize}
  \item \textbf{Planning}: manpower needs, policies, programmes.
  \item \textbf{Organizing}: structure of the HR department, tasks.
  \item \textbf{Directing}: motivating, guiding employees.
  \item \textbf{Coordinating} and \textbf{controlling}: audit of HR results against plan.
\end{itemize}
\lead{Role in the organization}
\begin{itemize}
  \item Supplies the \textbf{right people} at the right time.
  \item Builds \textbf{skills} (training) $\rightarrow$ productivity.
  \item \textbf{Fair pay}, incentives $\rightarrow$ motivation, low turnover.
  \item \textbf{Industrial peace}: handles grievances, unions.
  \item \textbf{Welfare and safety} $\rightarrow$ morale, fewer accidents.
  \item Advises line managers; keeps records; ensures labour-law compliance
    (Labour Act 2074 in Nepal).
  \item Aligns people with strategy $\rightarrow$ growth and development of the firm.
\end{itemize}}{%
\lead{Operative functions}
\begin{enumerate}
  \item \textbf{Procurement}: manpower planning, job analysis, recruitment, selection,
    placement, induction.
  \item \textbf{Development}: training, executive development, career planning, promotion.
  \item \textbf{Compensation}: job evaluation, wage and salary administration, incentives.
  \item \textbf{Integration}: reconcile employee and organizational interests: communication,
    grievance handling, collective bargaining.
  \item \textbf{Maintenance}: health, safety, welfare, social security, ergonomics.
  \item \textbf{Motivation}: monetary and non-monetary incentives.
  \item \textbf{Records and separation}: personnel records, transfers, retirement, exit.
\end{enumerate}
\textbf{Principles}: scientific selection, fair reward, high morale, dignity of labour,
team spirit, effective communication, joint management, effective utilization.}

\creamq{How does human resource management differ from personnel management?}{\m{2}}
{\color{sub}\footnotesize \tP{1} \yr{\textbf{71 Ch}} \gap \added}
\par\vspace{2pt}
\begin{tabularx}{\textwidth}{@{}L{2.2cm}XX@{}}
\toprule
\hrow \thead{Basis} & \thead{Personnel management} & \thead{Human resource management} \\
\midrule
View of people & Cost, a tool & Asset, a resource to develop \\
Nature & Routine, administrative, reactive & Strategic, proactive \\
Focus & Rules, records, compliance, pay & Development, motivation, culture, performance \\
Job design & Division of labour & Teamwork, empowerment \\
Responsibility & Personnel department only & All line managers share it \\
\bottomrule
\end{tabularx}

\penalty0\vspace{2pt}

%=======================================================================
\T{2.2 Personnel Policy}

\Q{Personnel policy: what it must include, its importance, and why it is discussed with new employees}{\m{3}\m{4}\m{5}\m{8}}
{\color{sub}\footnotesize \tF{7} \yr{\textbf{82 Bh}, 80 Ba, \textbf{79 Bh}, 75 Ash, \textbf{74 Ch}, 74 Ash, 73 Shr} \gap$\cdot$\gap 79 Bh asks it twice (policies; discussing at hiring)}

\begin{itemize}
  \item \textbf{Personnel policy} (\textbf{employee handbook}): written statement of what the
    employer expects from employees and what they may expect in return. A \textbf{guide for
    quick, consistent decisions}; a means to an end, not an objective.
  \item Eg a ``promotion from within'' policy serves the objective of employee development.
  \item Objectives: maximum development and use of workers' abilities; good labour relations;
    protect interests of labour and customers.
\end{itemize}

\sbs{%
\lead{Terms a good personnel policy includes}
\begin{enumerate}
  \item Welcome message; company name, history, values, objectives, structure.
  \item \textbf{Recruitment and selection}: sources, procedure, probation.
  \item \textbf{Working hours}, attendance, leave (sick, annual, holidays).
  \item \textbf{Pay}: salary structure, allowances, overtime, incentives, bonus.
  \item \textbf{Benefits}: insurance, provident fund, gratuity, pension, education aid.
  \item \textbf{Training} and development; \textbf{promotion and transfer} policy.
  \item \textbf{Code of conduct}: discipline, dress, harassment, conflict of interest.
  \item \textbf{Health and safety}; accident procedure.
  \item \textbf{Grievance procedure}; union and collective bargaining.
  \item \textbf{Separation}: resignation notice, retirement, dismissal, exit process.
\end{enumerate}}{%
\lead{Importance}
\begin{itemize}
  \item \textbf{Consistency and fairness}: same rule for everyone $\rightarrow$ trust.
  \item Quick decisions by managers without referring upward.
  \item Reduces conflict and grievances; legal protection for both sides.
  \item Attracts and retains good staff; builds culture.
\end{itemize}
\lead{Why discuss it at hiring / with new employees}
\begin{itemize}
  \item Employee knows \textbf{rights and duties} from day one: pay, leave, benefits,
    discipline.
  \item \textbf{No surprises} later $\rightarrow$ fewer disputes, lower early turnover.
  \item Part of \textbf{induction}: faster adjustment, sense of belonging.
  \item Written acknowledgement protects the employer legally.
  \item Clarifies expectations on conduct, safety, harassment.
\end{itemize}}

\penalty0\vspace{2pt}

%=======================================================================
\T{2.3 Manpower Planning}

\Q{Manpower (human resource) planning: meaning, steps, and why it is needed}{\m{3}\m{4}\m{8}}
{\color{sub}\footnotesize \tS{8} \yr{\textbf{81 Bh}, 81 Ba, 80 Ba, \textbf{79 Bh}, \textbf{76 Ch}, 76 Ash, \textbf{75 Ch}, 73 Shr} \gap$\cdot$\gap short note in 75 Ch}

\sbs{%
\begin{itemize}
  \item \textbf{Definition}: the process of having the \textbf{right number} of the
    \textbf{right kind} of people at the \textbf{right place} and \textbf{right time}, doing
    the jobs they suit.
  \item \textbf{Manpower} = total skills, knowledge, talent, creativity and values of the
    workforce.
  \item Two objectives: \textbf{use present employees fully}; \textbf{fill future needs}.
  \item Types: macro (national) vs micro (firm); short, medium, long term.
  \item Factors: working hours, shifts, nature of production, product mix, performance rate,
    hours lost (absenteeism), turnover.
\end{itemize}
\lead{Steps of the HR planning process}
\begin{enumerate}
  \item Study organizational objectives and plans.
  \item \textbf{Inventory of present manpower} (skills, age, department).
  \item \textbf{Forecast demand} (workload, expansion, technology, budget).
  \item \textbf{Forecast supply}: internal (promotion, transfer, retirement, turnover) and
    external (labour market).
  \item \textbf{Net requirement} = demand $-$ supply: shortage or surplus.
  \item Action plans: recruitment, training, redeployment, or downsizing.
  \item Monitor and review.
\end{enumerate}}{%
\lead{Importance / need}
\begin{itemize}
  \item \textbf{No shortage, no surplus}: avoids idle labour and overload.
  \item \textbf{Cuts labour cost}; right people raise productivity.
  \item Anticipates \textbf{retirement, turnover, expansion}: vacancies filled in time.
  \item Basis for \textbf{recruitment, training, career planning}.
  \item Copes with \textbf{technological change}: new skills planned ahead.
  \item Attracts talent; supports succession planning.
  \item Nepal: high outmigration of skilled youth $\rightarrow$ planning replacements and
    retention is essential.
\end{itemize}}

\penalty0\vspace{2pt}

%=======================================================================
\T{2.4 Recruitment and Selection}

\Q{Explain the process of recruitment and selection of manpower (steps of hiring); scientific selection}{\m{3}\m{5}\m{6}\m{8}}
{\color{sub}\footnotesize \tS{13} \yr{82 Ba, 81 Ba, \textbf{80 Bh}, \textbf{79 Bh}, \textbf{78 Bh}, \textbf{72 Ch}, \textbf{71 Ch}, 71 Shr, \textbf{69 Ch}, \textbf{68 Ch}, \textbf{67 Asa}, \textbf{66 Bh}, \textbf{65 Shr}}}

\begin{itemize}
  \item \textbf{Recruitment}: \textbf{searching} for prospective employees and
    \textbf{stimulating them to apply}; creates a pool (positive process).
  \item \textbf{Selection}: \textbf{choosing} the most suitable from the pool by
    eliminating the rest (negative process).
  \item \textbf{Hiring}: the final act of appointing, with offer, contract, placement.
\end{itemize}

\sbs{%
\lead{Sources of recruitment}
\begin{itemize}
  \item \textbf{Internal}: promotion, transfer, internal posting, former employees,
    employee referrals.
  \item \textbf{External}: advertisement (newspaper, websites like merojob, LinkedIn),
    campus recruitment, employment agencies, walk-ins, unions, consultants.
\end{itemize}
\lead{Steps of recruitment and selection}
\begin{enumerate}
  \item \textbf{Job analysis}, job description, staff requisition.
  \item \textbf{Advertise} the vacancy (title, duties, qualification, deadline).
  \item \textbf{Receive applications}; screen (shortlist).
  \item \textbf{Preliminary interview}: eliminate unfit early.
  \item \textbf{Application blank} (detailed form).
  \item \textbf{Employment tests}: written, aptitude, intelligence, trade, personality.
  \item \textbf{Employment interview}: HR, then department head.
  \item \textbf{Reference and background check}.
  \item \textbf{Medical examination}.
  \item \textbf{Final selection}, offer letter, appointment.
  \item \textbf{Placement} and \textbf{induction} (orientation, handbook).
\end{enumerate}}{%
\lead{Scientific selection}
\begin{itemize}
  \item Taylor's principle: \textbf{right person for the right job}, one who maintains the
    output with least effort.
  \item Selection by \textbf{objective tests and standards}, not favour or guess.
  \item \textbf{Fitness factors} checked:
  \begin{itemize}
    \item Physical: health, eyesight, strength.
    \item Personal: age, experience, background.
    \item Proficiency: present skill and ability.
    \item Competency: capacity to learn.
    \item Temperament, character, vocational interest.
  \end{itemize}
  \item \textbf{Methods}: psychological and aptitude tests, trade tests, structured
    interviews, work samples, assessment centres, medical tests.
\end{itemize}
\lead{Recruitment vs selection vs hiring}
{\small
\begin{tabularx}{\linewidth}{@{}L{1.5cm}XX@{}}
\toprule
\hrow \thead{Basis} & \thead{Recruitment} & \thead{Selection / hiring} \\
\midrule
Aim & Attract many applicants & Pick the best, appoint \\
Nature & Positive & Negative (eliminating) \\
Order & First & After recruitment \\
Output & Pool of candidates & Employee with contract \\
\bottomrule
\end{tabularx}}}

\creamq{Outsourcing}{\m{2}\m{3}}
{\color{sub}\footnotesize \tP{2} \yr{\textbf{80 Bh}, \textbf{69 Ch}} \gap \added}
\begin{itemize}
  \item Contracting a function to an \textbf{outside firm} instead of hiring staff for it.
    Eg security, cleaning, payroll, IT support, call centres; Nepal's IT firms do outsourced
    work for foreign clients.
  \item In recruitment: an agency does the hiring, or staff are supplied by a manpower
    company.
  \item $+$ lower cost, specialist skill, focus on core work, flexible headcount.
  \item $-$ less control, quality and confidentiality risk, dependence, staff morale.
\end{itemize}

\creamq{Referral recruitment: pros and cons}{\m{4}}
{\color{sub}\footnotesize \tP{1} \yr{\textbf{74 Ch}} \gap \added}
\begin{itemize}
  \item Existing employees recommend candidates.
  \item $+$ fast, cheap; pre-screened, trusted candidates; better cultural fit; referred
    staff stay longer; referral bonus motivates staff.
  \item $-$ nepotism and favouritism; less diversity; conflict if the referral fails or is
    rejected; closes the door to better outside talent.
  \item Use with the same tests and interviews as other candidates.
\end{itemize}

\creamq{Role of interview in hiring; the interviewing process}{\m{3}}
{\color{sub}\footnotesize \tP{2} \yr{\textbf{75 Ch}, \textbf{73 Ch}}}
\sbs{%
\begin{itemize}
  \item \textbf{Role}: get the maximum information about the candidate's suitability that
    tests cannot show: communication, attitude, confidence, motivation, fit with team.
  \item Two-way: candidate learns about the job and firm.
  \item Types: preliminary, structured, unstructured, panel, stress, technical, online.
\end{itemize}}{%
\begin{enumerate}
  \item \textbf{Preparation}: study job description and application; set questions.
  \item \textbf{Reception}: put candidate at ease.
  \item \textbf{Information exchange}: questions on knowledge, experience, attitude.
  \item \textbf{Close}: candidate's questions, next steps.
  \item \textbf{Evaluation}: rate on a standard form.
\end{enumerate}}

\penalty0\vspace{2pt}

%=======================================================================
\T{2.5 Training and Development}

\Q{Training and development: meaning, need, methods, and how training needs are identified}{\m{2}\m{3}\m{4}\m{5}}
{\color{sub}\footnotesize \tF{5} \yr{82 Ba, \textbf{81 Bh}, \textbf{79 Bh}, \textbf{72 Ch}, 71 Shr}}

\begin{itemize}
  \item \textbf{Training}: planned learning of \textbf{specific knowledge and skills} to do
    the \textbf{present job} better; short term; mainly for workers.
  \item \textbf{Development}: long-term education for \textbf{future roles}, conceptual and
    managerial growth; mainly for executives.
  \item Objectives: knowledge, skills, attitude change.
\end{itemize}

\sbs{%
\lead{Need / role (esp.\ for new employees)}
\begin{itemize}
  \item New hires become productive in the \textbf{minimum time}; no learning by mistakes.
  \item \textbf{Higher productivity}, quality; less wastage and spoilage.
  \item \textbf{Fewer accidents}, less machine breakdown.
  \item \textbf{Less supervision} needed.
  \item \textbf{Morale, job satisfaction}; lower turnover and absenteeism.
  \item Keeps pace with \textbf{new technology}; flexibility.
  \item Career growth, higher earnings; builds future managers.
\end{itemize}
\lead{Identifying training needs}
\begin{itemize}
  \item \textbf{Organization analysis}: goals, problems (low output, high scrap, accidents).
  \item \textbf{Task analysis}: job description vs skills required.
  \item \textbf{Person analysis}: performance appraisal gaps, supervisor reports,
    self-assessment.
  \item Other signals: complaints, new machines, promotions, surveys, exit interviews.
\end{itemize}}{%
\lead{Methods: workers}
\begin{itemize}
  \item \textbf{On-the-job}: supervisor trains at the workplace; most effective, realistic.
  \item \textbf{Demonstration}: instructor shows the task (Eg lathe operation).
  \item \textbf{Vestibule}: separate school duplicating real equipment; many trainees, no
    production pressure; costly.
  \item \textbf{Apprenticeship}: long, for all-round craftsmen (fitters, welders).
  \item \textbf{Special courses}: knowledge-based, in-house or outside institutes.
\end{itemize}
\lead{Methods: supervisors and executives}
\begin{itemize}
  \item Induction and orientation; lectures; conferences.
  \item \textbf{On the job}: understudy, job rotation, committee membership, job
    enlargement / enrichment, MBO.
  \item \textbf{Off the job}: case study, business games, role playing, seminars,
    e-learning.
\end{itemize}}

\creamq{HR problems: talent poaching; short-term employee development programme; appropriate vs inappropriate human resources}{\m{4}\m{8}}
{\color{sub}\footnotesize \tP{2} \yr{79 Ba, 73 Shr} \gap \added}
\sbs{%
\lead{Tackling talent poaching (79 Ba)}
\begin{itemize}
  \item Competitive pay, reviewed against the market.
  \item Retention bonus, ESOP, bond or notice clauses.
  \item Clear career path, promotion from within.
  \item Training, recognition, good culture, flexible work.
  \item Stay interviews; fix reasons found in exit interviews.
  \item Succession planning so no one person is critical.
\end{itemize}
\lead{Short-term development programme}
\begin{enumerate}
  \item Identify skill gaps (appraisal).
  \item Set 3--6 month objectives per employee.
  \item Workshops, mentoring, job rotation, online courses.
  \item Review, reward, plan the next cycle.
\end{enumerate}}{%
\lead{Appropriate vs inappropriate human resources (73 Shr)}
\begin{itemize}
  \item \textbf{Appropriate}: people whose skills, qualification, attitude and number
    \textbf{match} job requirements $\rightarrow$ productive, satisfied.
  \item \textbf{Inappropriate}: mismatch: over-qualified (bored, leave), under-qualified
    (errors, need supervision), surplus (idle cost), wrong attitude.
  \item Ideas to get appropriate HR: manpower planning, job analysis, scientific selection,
    training, appraisal, redeployment, fair pay.
\end{itemize}}

\penalty0\vspace{2pt}

%=======================================================================
\T{2.6 Job Analysis, Job Evaluation and Merit Rating}

\Q{Job analysis: meaning, methods, job description and job specification}{\m{2}\m{3}\m{4}\m{5}\m{6}}
{\color{sub}\footnotesize \tS{10} \yr{\textbf{82 Bh}, \textbf{81 Bh}, 81 Ba, \textbf{79 Bh}, \textbf{78 Bh}, \textbf{71 Ch}, \textbf{68 Ch}, \textbf{67 Asa}, \textbf{66 Bh}, \textbf{65 Shr}} \gap$\cdot$\gap job design as a short note in 67 Asa}

\begin{itemize}
  \item \textbf{Job}: group of tasks, duties and responsibilities done by one person for pay.
  \item \textbf{Job analysis}: systematic study of a job to find its \textbf{tasks} and the
    \textbf{skills, knowledge, abilities and responsibilities} needed to do it well. Answers
    what, how, why, with what tools, under what conditions.
  \item Outputs: \textbf{job description} (about the job) and \textbf{job specification}
    (about the person). Used for recruitment, training, job evaluation, pay, appraisal.
\end{itemize}

\sbs{%
\lead{Steps}
\begin{enumerate}
  \item Collect background info (charts, old descriptions).
  \item Select jobs to analyse.
  \item \textbf{Collect job data} (methods below).
  \item Process and verify with job holder and supervisor.
  \item Prepare job description.
  \item Prepare job specification.
\end{enumerate}
\lead{Methods of collecting data}
\begin{itemize}
  \item \textbf{Questionnaire}: holders fill structured forms; cheap for many jobs;
    misunderstanding possible.
  \item \textbf{Interview}: holder describes duties, even unobservable ones; may exaggerate.
  \item \textbf{Observation}: analyst watches the work; firsthand; best for manual jobs;
    workers may change behaviour.
  \item \textbf{Diary / log}, \textbf{critical incident}, \textbf{technical conference} with
    experts.
\end{itemize}
\lead{Most suitable in Nepal? (82 Bh)}
\begin{itemize}
  \item \textbf{Interview + observation}: many small firms have no written records, jobs are
    manual or mixed, literacy varies $\rightarrow$ questionnaires return poor data.
    Large firms and banks can add questionnaires.
\end{itemize}}{%
\lead{Job description (what the job is)}
\begin{itemize}
  \item Job title, code, department, location.
  \item Job summary (purpose).
  \item Duties and responsibilities: what, how, why.
  \item Reporting relationships.
  \item Machines, tools, materials.
  \item Working conditions: shifts, hazards.
  \item Pay grade, benefits.
\end{itemize}
\lead{Job specification (who can do it)}
\begin{itemize}
  \item Education, training, experience.
  \item Skills, knowledge, abilities.
  \item Physical and mental requirements.
  \item Personal qualities (communication, attitude).
\end{itemize}
\lead{Example: Lecturer, engineering college (81 Ba)}
{\small
\begin{itemize}
  \item \textbf{JD}: title Lecturer, Dept of Electronics; reports to HoD. Teach 16--18
    periods/week, prepare lessons, lab supervision, set and evaluate exams, guide projects,
    research and publish, departmental duties.
  \item \textbf{JS}: Master's in relevant engineering, first division; NEC registration;
    teaching or industry experience preferred; communication and presentation skill,
    lab skills, ethics, research aptitude.
\end{itemize}}
\lead{Job design (short note)}
\begin{itemize}
  \item Structuring tasks of a job for productivity and satisfaction: job simplification,
    rotation, enlargement (more tasks), enrichment (more responsibility), ergonomics.
\end{itemize}}

\Q{Job evaluation: meaning, methods, and comparison with merit rating}{\m{2}\m{3}}
{\color{sub}\footnotesize \tP{2} \yr{81 Ba, \textbf{79 Bh}}}
\sbs{%
\begin{itemize}
  \item Systematic determination of the \textbf{relative worth of a job} compared with
    others $\rightarrow$ fair, uniform \textbf{wage structure}.
  \item Principle: \textbf{rate the job, not the man}.
  \item \textbf{Non-quantitative methods}:
  \begin{itemize}
    \item \textbf{Ranking}: jobs ranked whole, most to least important.
    \item \textbf{Classification / grading}: jobs slotted into predefined grades.
  \end{itemize}
  \item \textbf{Quantitative methods}:
  \begin{itemize}
    \item \textbf{Factor comparison}: skill, mental effort, physical effort,
      responsibility, working conditions compared factor by factor with key jobs.
    \item \textbf{Point rating}: points per factor degree, totals set the grade.
  \end{itemize}
\end{itemize}}{%
{\small
\begin{tabularx}{\linewidth}{@{}L{1.5cm}XX@{}}
\toprule
\hrow \thead{Basis} & \thead{Job evaluation} & \thead{Merit rating} \\
\midrule
Rates & The job & The person doing it \\
Purpose & Wage structure for the job & Reward, promotion over the base wage \\
Timing & Before hiring, on job change & Periodically, after hiring \\
Basis & Job analysis & Actual performance, traits \\
Output & Job grades & Ranking of employees \\
\bottomrule
\end{tabularx}}}

\Q{Merit rating / performance appraisal: meaning, need, traditional and modern methods}{\m{2}\m{4}\m{5}\m{6}}
{\color{sub}\footnotesize \tF{7} \yr{82 Ba, \textbf{81 Bh}, \textbf{79 Bh}, \textbf{78 Bh}, \textbf{70 Ch}, \textbf{68 Ba}, \textbf{66 Bh}}}
\sbs{%
\begin{itemize}
  \item \textbf{Merit rating}: systematic, periodic assessment of an \textbf{employee's
    worth} in terms of job performance, integrity, leadership, behaviour, aptitude.
  \item Factors: quality and quantity of work, job knowledge, dependability, attitude,
    initiative, versatility, leadership.
\end{itemize}
\lead{Why necessary}
\begin{itemize}
  \item Basis for \textbf{promotion, transfer, pay rise, bonus}.
  \item Feedback to employee $\rightarrow$ improvement.
  \item Identifies \textbf{training needs}.
  \item Validates selection; objective, fair decisions; motivation.
\end{itemize}
\lead{Traditional methods (trait-based)}
\begin{itemize}
  \item \textbf{Ranking}: best to worst, overall.
  \item \textbf{Paired / person-to-person comparison}: each with every other.
  \item \textbf{Man-to-man comparison}: against key men on 5 traits, 5 degrees each.
  \item \textbf{Grading}: predefined grades (outstanding, satisfactory, unsatisfactory).
  \item \textbf{Graphic rating scale}: traits rated on a scale; most common.
  \item \textbf{Checklist}: yes/no statements.
  \item \textbf{Free essay}: supervisor's written description.
  \item \textbf{Critical incident}: record of good and bad events.
\end{itemize}}{%
\lead{Modern methods (result-based)}
\begin{itemize}
  \item \textbf{MBO / appraisal by results}: manager and employee jointly set goals, appraise
    against them.
  \item \textbf{Assessment centre}: simulations, group exercises, in-basket tasks observed by
    trained assessors; for promotion.
  \item \textbf{360-degree feedback}: from superiors, peers, subordinates, customers, self.
    \added
  \item \textbf{BARS} (behaviourally anchored rating scale): scale points tied to real
    behaviours. \added
  \item \textbf{Human resource accounting}: employee's value in money. \added
  \item \textbf{KPI / balanced scorecard} dashboards. \added
\end{itemize}
\lead{Merits of modern methods}
\begin{itemize}
  \item Objective, measurable; less bias.
  \item Employee participates $\rightarrow$ commitment.
  \item Links appraisal to goals and development.
\end{itemize}}

\penalty0\vspace{2pt}

%=======================================================================
\T{2.7 Wages and Incentives}

\Q{Wages and salary: difference, factors affecting the wage/salary structure, how wages are calculated, salary and benefits you expect}{\m{2}\m{3}\m{4}\m{5}\m{8}}
{\color{sub}\footnotesize \tS{12} \yr{\textbf{82 Bh}, 82 Ba, \textbf{80 Bh}, \textbf{78 Bh}, \textbf{76 Ch}, \textbf{75 Ch}, 74 Ash, \textbf{73 Ch}, \textbf{72 Ch}, 72 Ka, 70 Asa, \textbf{68 Ba}}}

\sbs{%
\begin{itemize}
  \item \textbf{Wage}: remuneration for labour, paid per hour, day, week or per unit of
    output; production and manual workers.
  \item \textbf{Salary}: fixed monthly or annual pay; office, professional, managerial
    staff.
  \item \textbf{Nominal (money) wage}: cash only. \textbf{Real wage}: what the cash buys, plus
    benefits.
  \item Levels: \textbf{minimum wage} (legal floor; Nepal fixes it periodically),
    \textbf{fair wage} (equal pay for equal work), \textbf{living wage} (decent family
    living).
  \item \textbf{Structure}: job classes / grades from job evaluation; flat rate or rate range
    per grade; \textbf{direct} pay (basic, overtime, bonus) + \textbf{indirect} (insurance,
    provident fund, gratuity, allowances).
\end{itemize}
\lead{Salary and benefits I would expect (82 Ba, 75 Ch)}
\begin{itemize}
  \item Basic salary at or above market rate for the post, with annual increment.
  \item Allowances: dearness, transport, communication, festival (Dashain) allowance.
  \item Provident fund, gratuity / SSF contribution; health and accident insurance.
  \item Paid leave; overtime; performance bonus.
  \item Training, career growth; flexible hours. Justify with qualification and market data.
\end{itemize}}{%
\lead{Factors affecting wage / salary structure}
\begin{enumerate}
  \item \textbf{Ability of the firm to pay} (profit, size).
  \item \textbf{Demand and supply} of labour.
  \item \textbf{Prevailing market rate} for comparable jobs.
  \item \textbf{Cost of living} (inflation index).
  \item \textbf{Legal provisions}: minimum wage, Labour Act, SSF.
  \item \textbf{Job requirements} from job evaluation: skill, effort, responsibility,
    hazards.
  \item \textbf{Productivity} (output per man-hour).
  \item \textbf{Bargaining power} of trade unions.
  \item Employee's qualification, experience, seniority.
  \item Regularity of employment, working hours, promotion prospects, training cost.
\end{enumerate}
\textbf{Examples} (80 Bh): an IT engineer in Kathmandu earns more than a same-degree
engineer in a small town (supply-demand and market rate); garment workers' pay rises when the
government revises the minimum wage (legal); a union-negotiated increase at NEA
(bargaining power).}

\Q{Incentives: meaning, need, types}{\m{2}\m{3}\m{4}}
{\color{sub}\footnotesize \tS{10} \yr{\textbf{78 Bh}, \textbf{76 Ch}, 74 Ash, \textbf{73 Ch}, \textbf{72 Ch}, 72 Ka, 71 Shr, 70 Asa, \textbf{68 Ba}, \textbf{66 Bh}} \gap$\cdot$\gap ``incentive programs'' short note in 68 Ba, 66 Bh}
\sbs{%
\begin{itemize}
  \item \textbf{Incentive}: payment or benefit \textbf{over and above} normal wage, linked to
    output or performance above a standard, to motivate higher productivity.
\end{itemize}
\lead{Why needed}
\begin{itemize}
  \item Raises productivity, lowers unit cost.
  \item Rewards efficient workers fairly.
  \item Motivation, morale, job satisfaction.
  \item Reduces supervision, absenteeism, turnover.
  \item Attracts and retains talent.
\end{itemize}}{%
\lead{Types}
\begin{itemize}
  \item \textbf{Financial}: bonus, profit sharing, commission, piece-rate premium, ESOP.
  \item \textbf{Non-financial}: praise, recognition, promotion, job security, training,
    flexible hours, better conditions.
  \item \textbf{Semi-financial}: subsidized lunch, medical, recreation, pension.
  \item \textbf{Individual}: based on one person's output (Taylor, Halsey, Rowan).
  \item \textbf{Group}: based on team or plant output (Scanlon, Priestman, profit sharing):
    efficient workers help slow ones.
\end{itemize}}

\creamq{How wages are calculated: wage plans and incentive formulas}{\m{3}}
{\color{sub}\footnotesize \tP{1} \yr{\textbf{75 Ch}} \gap$\cdot$\gap no numerical has been set in 29 papers; each formula shown once}

\sbs{%
\lead{Basic plans}
\begin{itemize}
  \item \textbf{Time rate}: wage $= T_a \times R$ (hours worked $\times$ hourly rate).
    Simple, secure; no push for output.
  \item \textbf{Piece rate}: wage $=$ units $\times$ rate per unit. Rewards output; no
    security.
  \item \textbf{Combined}: guaranteed time wage + piece earnings above it.
\end{itemize}
\lead{Taylor's differential piece rate}
\begin{itemize}
  \item Two rates: low (Eg 80\%) below standard, high (Eg 120\%) at or above.
  \item Standard 80 jobs/day (8 h), Rs 12/h $\rightarrow$ standard rate $12/10 =$ Rs 1.20.
    A makes 75: $75 \times 0.96 =$ \textbf{Rs 72}. B makes 85: $85 \times 1.44 =$
    \textbf{Rs 122.40}.
\end{itemize}
\lead{Halsey premium plan (1891)}
\begin{itemize}
  \item Bonus $= \tfrac12 R\,(T_s - T_a)$; wage $= T_a R + \tfrac12 R\,(T_s - T_a)$.
  \item $T_s = 8$, $T_a = 6$, $R =$ Rs 10: bonus Rs 10, wage \textbf{Rs 70}.
\end{itemize}
\lead{Rowan plan (1901)}
\begin{itemize}
  \item Bonus $= \dfrac{T_s - T_a}{T_s}\,T_a R$. Peaks when $T_a = T_s/2$.
\end{itemize}
{\small
\begin{tabularx}{\linewidth}{@{}XXXXXXXX@{}}
\toprule
\hrow $T_a$ & 7 & 6 & 5 & 4 & 3 & 2 & 1 \\
Bonus & 8.75 & 15.00 & 18.75 & \textbf{20.00} & 18.75 & 15.00 & 8.75 \\
\bottomrule
\end{tabularx}}
{\color{sub}\scriptsize $T_s = 8$ h, $R =$ Rs 10; bonus in Rs. Computed here (source table p184 is
blank in the scan).}}{%
\lead{Emerson efficiency plan}
\begin{itemize}
  \item Efficiency $=$ actual / standard output $\times 100$.
  \item Below 66.67\%: time wage only; 66.67--100\%: rising bonus to 20\% at 100\%; above
    100\%: 20\% + 1\% per 1\% extra.
  \item Rs 40/day, standard 50 pins: X 30 pins (60\%) Rs 40; Y 50 (100\%) Rs 48; Z 60
    (120\%) \textbf{Rs 56}.
\end{itemize}
\lead{Scanlon plan (group)}
\begin{itemize}
  \item Bonus \% $= \dfrac{L_b/S_b - L_a/S_a}{L_b/S_b} \times 100$ (labour cost / sales).
  \item Base 1.35 / 12.40 = 10.887\%; month 1.45 / 15.50 = 9.355\% $\rightarrow$ bonus
    \textbf{14.07\%}. With 2.15 lakh labour cost: 13.87\% $>$ base, no bonus.
\end{itemize}
\lead{Priestman plan (group)}
\begin{itemize}
  \item Bonus \% $= \dfrac{P_a/W_a - P_b/W_b}{P_b/W_b} \times 100$ (output per worker).
  \item Standard 1200 cars / 500 workers = 2.4. Next month 1400 cars: same 500 workers
    $\rightarrow$ 2.8, bonus \textbf{16.67\%}; 600 workers $\rightarrow$ 2.333, $-2.78\%$,
    no bonus.
\end{itemize}
\lead{Profit sharing}
\begin{itemize}
  \item Agreed \% of annual profit shared with employees (Nepal's Labour Act and Bonus Act
    require a share of profit as bonus).
\end{itemize}}

\penalty0\vspace{2pt}

%=======================================================================
\T{2.8 Industrial Relations and Collective Bargaining}

\Q{Industrial relations; collective bargaining and its process}{\m{3}\m{4}\m{6}}
{\color{sub}\footnotesize \tF{4} \yr{\textbf{68 Ch}, \textbf{67 Asa}, \textbf{66 Bh}, \textbf{65 Shr}} \gap$\cdot$\gap 2065--2068 papers only \gap \added}
\sbs{%
\lead{Industrial relations}
\begin{itemize}
  \item Relations between \textbf{management, workers (and their unions) and government}
    in an industry.
  \item Aim: \textbf{industrial peace}, higher productivity, fair treatment, no strikes or
    lockouts.
  \item Tools: collective bargaining, grievance procedure, joint consultation, workers'
    participation, conciliation, arbitration, labour courts.
  \item Causes of poor relations: low wages, bad conditions, unfair dismissal, political
    unions, poor communication.
  \item Nepal: Labour Act 2074, Trade Union Act 2049; frequent politically affiliated
    strikes.
\end{itemize}}{%
\lead{Collective bargaining}
\begin{itemize}
  \item \textbf{Negotiation between employer and a union representing workers} on wages,
    hours, conditions, benefits; result is a written \textbf{collective agreement}.
\end{itemize}
\lead{Process}
\begin{enumerate}
  \item \textbf{Preparation}: union frames demands (charter); management studies data.
  \item \textbf{Negotiation}: meetings, proposals, counter-proposals, give and take.
  \item \textbf{Agreement}: terms written, signed, registered with labour office.
  \item \textbf{Implementation} of the agreement.
  \item \textbf{Administration and review}: grievances handled; renegotiated at expiry.
  \item Deadlock $\rightarrow$ conciliation, mediation, arbitration.
\end{enumerate}}

\penalty0\vspace{2pt}

%=======================================================================
\T{2.9 PYQ Mapping}

\lead{Theory questions, covered in this document}
\begin{itemize}
  \item \textbf{Definition, functions, role of personnel management} $\rightarrow$ $\S$2.1.
    17 papers; HRM vs personnel management \yr{\textbf{71 Ch} \m{2}}.
  \item \textbf{Personnel policy / employee handbook} $\rightarrow$ $\S$2.2. 7 papers,
    \m{3} to \m{5}; 75 Ash prints no marks against it.
  \item \textbf{Manpower / HR planning} $\rightarrow$ $\S$2.3. 8 papers.
  \item \textbf{Recruitment and selection, steps of hiring, scientific selection,
    recruitment vs selection / hiring} $\rightarrow$ $\S$2.4. 13 papers.
  \item \textbf{Outsourcing} \yr{\textbf{80 Bh} \m{3}, \textbf{69 Ch} \m{2}}; \textbf{referral recruitment}
    \yr{\textbf{74 Ch} \m{4}}; \textbf{interview} \yr{\textbf{75 Ch} \m{3}, \textbf{73 Ch} \m{3}} $\rightarrow$ $\S$2.4.
  \item \textbf{Training: need, types, methods, training needs} $\rightarrow$ $\S$2.5.
    \yr{82 Ba \m{4}, \textbf{81 Bh} \m{3+2}, \textbf{79 Bh} \m{3}, \textbf{72 Ch} \m{4}, 71 Shr \m{5}}
  \item \textbf{Talent poaching, development programme, appropriate HR} \yr{79 Ba \m{4+4},
    73 Shr \m{8}} $\rightarrow$ $\S$2.5.
  \item \textbf{Job analysis, job description, job specification, lecturer example, method
    for Nepal, job design} $\rightarrow$ $\S$2.6. 10 papers.
  \item \textbf{Job evaluation vs merit rating} \yr{81 Ba \m{3}, \textbf{79 Bh} \m{2}};
    \textbf{merit rating / performance appraisal methods} (7 papers) $\rightarrow$ $\S$2.6.
  \item \textbf{Wage vs salary, factors of wage structure, how wages are calculated, expected
    salary} $\rightarrow$ $\S$2.7. 12 papers.
  \item \textbf{Incentives, why needed, incentive programs} $\rightarrow$ $\S$2.7. 10 papers.
  \item \textbf{Industrial relations, collective bargaining} $\rightarrow$ $\S$2.8.
    \yr{\textbf{68 Ch} \m{4}, \textbf{67 Asa} \m{4+6}, \textbf{66 Bh} \m{6}, \textbf{65 Shr} \m{4+3}}
\end{itemize}

\lead{Short-note prompts from the Detailed PYQ's ``Extra'' section, placed here}
\begin{itemize}
  \item Industrial relation \yr{\textbf{67 Asa}, \textbf{65 Shr}}; collective bargaining \yr{\textbf{68 Ch}}; job
    design and work efficiency \yr{\textbf{67 Asa}}; incentive programs \yr{\textbf{68 Ba}, \textbf{66 Bh}}; manpower
    planning \yr{\textbf{75 Ch}}.
  \item The Detailed PYQ files 66 Bh's ``What do you mean by industrial relation? [6]'' under
    Ch.\,1 (organizational structure); it is answered here.
\end{itemize}

\lead{Numerical content}
\begin{itemize}
  \item \textbf{No paper in 29 sets a wage or incentive calculation.} The formulas in $\S$2.7
    answer the 3-mark ``how are wages calculated?'' (75 Ch) and are worth knowing in case a
    numerical appears.
\end{itemize}

\lead{Corrections to the source notes}
\begin{itemize}
  \item Adhikari p190, Priestman example (b): $(1400/600 - 1200/500)/(1200/500) = -2.78\%$,
    not $-2.91\%$. The conclusion (no bonus) stands.
  \item Adhikari p181 spells the plan ``Hasley''; it is the \textbf{Halsey} plan (F.A.\ Halsey).
\end{itemize}
